%% 4-3 にた形と比：相似（第13回：組 V4-2。ページ制限なし。偶数ページ）
%% ほかの巻が指している形：
%%   5-10「断面の半径は、頂点からの長さ y に比例します（相似、4-3）」「頂点から y の所の断面は、1辺 a/h·y の正方形です（4-3）」
%% この章で定義した部品（頭文字 bkivd）：
%%   \bkivdstep{番号}{見出し}  説明の段の見出し（丸数字＋太字の1行）
%%   \bkivdline{式}            式を、まん中に1行で置く
%%   \bkivdmidashi{見出し}     本文の中の小見出し（それだけで1行）
%%   \bkivdright{(頂点)}{1本めの向きの角}{2本めの向きの角}  直角の印（2mm）
%%   \bkivdsun{(x,y)}{半径}    太陽（円と、まわりの短い線）
\makeatletter
\newcommand\bkivdstep[2]{\par\addvspace{6pt}\Needspace{5\baselineskip}\noindent
  {\bkcircled[bkOrange]{#1}\enspace\sffamily\bfseries\color{bkOrange!75!black}#2}\par\nobreak\vspace{2pt}\noindent\ignorespaces}
\newcommand\bkivdline[1]{\par\addvspace{3pt}{\centering #1\par}\addvspace{3pt}\noindent\ignorespaces}
\newcommand\bkivdmidashi[1]{\par\addvspace{4pt}\Needspace{6\baselineskip}\noindent\textbf{#1}\par\noindent\ignorespaces}
\newcommand\bkivdright[3]{\draw[bkInk, line width=0.5pt] ($#1+(#2:2.2)$) -- ($#1+(#2:2.2)+(#3:2.2)$) -- ($#1+(#3:2.2)$);}
\newcommand\bkivdsun[2]{\begin{scope}[shift={#1}]
  \fill[bkYellow] (0,0) circle (#2);
  \foreach \bk@a in {0,45,...,315} {\draw[bkYellow!80!black, line width=0.8pt] (\bk@a:#2*1.35) -- (\bk@a:#2*1.8);}\end{scope}}
\newcommand\bkivdpoint[1]{\par\addvspace{4pt}\noindent{\sffamily\bfseries\color{bkBlue!80!black}ポイント}\enspace #1\par\addvspace{4pt}}
\makeatother

\bkchapter{4}{3}{にた形と比：相似}{\bklevelrange{2}{3}}{にた形と比：相似}
\bkmichishirube{2-5,2-15,3-5,4-1,4-2}{4-4,4-6,5-10}

\begin{goal}
\begin{itemize}[label={\color{\bkTT}●}, leftmargin=1.6\zw, itemsep=0pt, topsep=1pt]
\item \bklevel{2}\enspace 拡大図・縮図と縮尺を使って、本当の長さや図の上の長さを求められる。
\item \bklevel{3}\enspace 三角形の相似条件を見つけ、相似比や平行線と線分の比から、辺の長さを求められる。
\item \bklevel{3}\enspace 面積の比は $k^2$、体積の比は $k^3$ になることを使い、すいの断面や直角三角形の辺の比を説明できる。
\end{itemize}
\end{goal}

\begin{hatena}
\noindent\begin{minipage}[c]{0.47\linewidth}
晴れた日に、校庭の木の高さを知りたくなりました。木はとても高くて、のぼることも、メジャーを当てることもできません。でも、地面にできた影の長さなら、メジャーではかれます。

木の影は\qty{9}{m}でした。同じ時刻に、長さ\qty{1}{m}の棒をまっすぐ立てると、影は\qty{1.5}{m}でした。木の高さは何mかな？
\end{minipage}\hfill
\begin{minipage}[c]{0.5\linewidth}\centering
\begin{tikzpicture}[x=0.9mm, y=0.9mm]
  % 1 m を 3.3 mm でかく。木の高さ 6 m（19.8 mm）、影 9 m（29.7 mm）。棒 1 m（3.3 mm）、影 1.5 m（4.95 mm）
  \fill[bkGray!20] (-14,-2.6) rectangle (50,0);
  \draw[bkGray, line width=0.8pt] (-14,0) -- (50,0);
  \bktree{(0,0)}{19.6}
  \fill[bkInk!55] (0,-0.8) rectangle (29.7,0);
  \draw[bkOrange, line width=0.7pt, dash pattern=on 2.5pt off 1.5pt] (-6,23.8) -- (32.7,-2);
  \draw[bkOrange, line width=0.7pt, dash pattern=on 2.5pt off 1.5pt] (30,8.6) -- (43.9,-0.67);
  \draw[bkInk, line width=1.4pt] (38,0) -- (38,3.3);
  \fill[bkInk!55] (38,-0.8) rectangle (42.95,0);
  \bkivdsun{(-9.5,28.5)}{2.4}
  \draw[bkInk, line width=0.5pt, {Stealth[length=2.5pt]}-{Stealth[length=2.5pt]}] (-9,0.3) -- (-9,19.5);
  \node[font=\sffamily\footnotesize, anchor=east] at (-9.8,10) {？};
  \node[font=\sffamily\footnotesize, anchor=north] at (14.85,-3.2) {影 \qty{9}{m}};
  \node[font=\sffamily\footnotesize, anchor=north] at (40.5,-3.2) {影 \qty{1.5}{m}};
  \node[font=\sffamily\footnotesize, anchor=south west] at (39,5) {棒 \qty{1}{m}};
\end{tikzpicture}
\end{minipage}
\end{hatena}

\noindent 木と棒は、大きさがまるでちがいます。でも、木とその影が作る三角形と、棒とその影が作る三角形は、「形」が同じです。この章では、形が同じで大きさだけがちがう図形を調べます。形が同じなら、長さの「比」が同じになります。その比を使うと、はかれない長さが計算で求められます。

%% ---------------- やってみよう・思い出そう ----------------
\begin{yattemiyou}[方眼で、2倍の拡大図をかこう]
\noindent\begin{minipage}[c]{0.5\linewidth}
\textbf{用意するもの}\par
方眼紙、定規、分度器
\begin{enumerate}[label={\bkcircled[bkGreen]{\arabic*}}, leftmargin=2.4\zw, itemsep=1pt, topsep=3pt]
\item 方眼紙に、右の図の三角形ABCをかこう。
\item 右へ何マス、上へ何マスを、すべて2倍にして、三角形DEFをかこう。
\item 2つの三角形の辺の長さと角の大きさを、定規と分度器ではかって、表にまとめよう。
\end{enumerate}
\end{minipage}\hfill
\begin{minipage}[c]{0.46\linewidth}\centering
\begin{tikzpicture}[x=0.86mm, y=0.86mm]
  % 1マス 4 mm。ABC：A(0,0) B(4,0) C(1,3)。DEF：D(6,0) E(14,0) F(8,6)（マス）
  \draw[bkLine, line width=0.3pt] (-2,-2) grid[step=4] (58,26);
  \fill[bkBlue!20] (0,0) -- (16,0) -- (4,12) -- cycle;
  \draw[bkBlue, line width=1.1pt] (0,0) -- (16,0) -- (4,12) -- cycle;
  \fill[bkOrange!20] (24,0) -- (56,0) -- (32,24) -- cycle;
  \draw[bkOrange, line width=1.1pt] (24,0) -- (56,0) -- (32,24) -- cycle;
  \foreach \p/\q/\l in {0/0/A, 16/0/B, 4/12/C, 24/0/D, 56/0/E, 32/24/F} {\fill[bkInk] (\p,\q) circle (0.8);}
  \node[font=\sffamily\footnotesize, anchor=north east, fill=white, inner sep=0.8pt] at (-0.4,-0.4) {A};
  \node[font=\sffamily\footnotesize, anchor=north, fill=white, inner sep=0.8pt] at (16,-0.6) {B};
  \node[font=\sffamily\footnotesize, anchor=south, fill=white, inner sep=0.8pt] at (4,12.8) {C};
  \node[font=\sffamily\footnotesize, anchor=north, fill=white, inner sep=0.8pt] at (24,-0.6) {D};
  \node[font=\sffamily\footnotesize, anchor=north west, fill=white, inner sep=0.8pt] at (56.4,-0.4) {E};
  \node[font=\sffamily\footnotesize, anchor=south, fill=white, inner sep=0.8pt] at (32,24.8) {F};
\end{tikzpicture}
\end{minipage}
\par\vspace{4pt}
ソラさんがはかった結果です（長さの単位は、方眼の1マス）。
\begin{bkcenter}
\begin{bktable}{colspec={Q[l,m,wd=24mm] *{6}{Q[c,m,wd=11mm]}}, rowsep=1.2pt, row{1}={font=\sffamily\small}, row{3}={font=\sffamily\small}}
 & AB & BC & CA & $\angle\mathrm{A}$ & $\angle\mathrm{B}$ & $\angle\mathrm{C}$ \\
三角形ABC & 4 & 4.2 & 3.2 & 72° & 45° & 63° \\
 & DE & EF & FD & $\angle\mathrm{D}$ & $\angle\mathrm{E}$ & $\angle\mathrm{F}$ \\
三角形DEF & 8 & 8.5 & 6.3 & 72° & 45° & 63° \\
\end{bktable}
\end{bkcenter}
\end{yattemiyou}

\noindent やってみようの答え：辺の長さは、DEがABの2倍、EFがBCの約2倍、FDがCAの約2倍です（はかったときのずれがあるので「約」）。角の大きさは、3つとも、もとの三角形と同じです。2倍の拡大図は、「長さはすべて2倍、角はそのまま」の図形です。

\begin{omoidasou}
\bkrecall{2-5}{比 $a:b$ の比の値は $a\div b$。比の両方に同じ数をかけても、わっても、等しい比になる。}
\bkrecall{2-15}{比例：$x$ が2倍、3倍になると、$y$ も2倍、3倍。$y=ax$。}
\bkrecall{3-5}{傾き＝たての変化量÷横の変化量。坂の急さを表す。}
\bkrecall{4-1}{ぴったり重なる図形は合同。三角形は「3辺」「2辺とその間の角」「1辺とその両はしの角」が決まると1つに決まる（合同条件）。平行線の同位角・錯角は等しい。対頂角は等しい。}
\end{omoidasou}

%% ---------------- 数学の道具 ----------------
\bksection{doguu}{数学の道具}

\bksubsection{1}{拡大図・縮図と縮尺}

\leavevmode\label{bk:4-3:kakudai}スマートフォンの写真を、2本の指で広げて拡大すると、写っている人の顔は大きくなります。でも、顔の形は変わりません。目と目の間も、鼻の長さも、口のはばも、同じ倍率で大きくなるからです。

\begin{bkfigure}
\begin{tikzpicture}[x=1mm, y=1mm]
  % 1マス 3.5 mm。もとの家（はば4・高さ4、屋根の上 6）
  \draw[bkLine, line width=0.3pt] (0,0) grid[step=3.5] (126,28);
  \begin{scope}[shift={(3.5,0)}]
    \fill[bkBlue!20] (0,0) -- (14,0) -- (14,7) -- (7,14) -- (0,7) -- cycle;
    \draw[bkBlue, line width=1.1pt] (0,0) -- (14,0) -- (14,7) -- (7,14) -- (0,7) -- cycle;
    \node[font=\sffamily\footnotesize, anchor=south] at (7,15.5) {もとの形};
  \end{scope}
  \begin{scope}[shift={(28,0)}]
    \fill[bkOrange!20] (0,0) -- (28,0) -- (28,14) -- (14,28) -- (0,14) -- cycle;
    \draw[bkOrange, line width=1.1pt] (0,0) -- (28,0) -- (28,14) -- (14,28) -- (0,14) -- cycle;
    \node[font=\sffamily\footnotesize, fill=white, inner sep=1pt] at (14,8) {2倍の拡大図};
  \end{scope}
  \begin{scope}[shift={(70,0)}]
    \fill[bkRed!15] (0,0) -- (42,0) -- (42,7) -- (21,14) -- (0,7) -- cycle;
    \draw[bkRed, line width=1.1pt] (0,0) -- (42,0) -- (42,7) -- (21,14) -- (0,7) -- cycle;
    \node[font=\sffamily\footnotesize, anchor=south, fill=white, inner sep=1pt] at (21,16) {横だけ3倍（形が変わる）};
  \end{scope}
\end{tikzpicture}
\bkfigcaption{たても横も2倍にすると、形は同じ（オレンジ）。横だけ3倍にすると、屋根の角が変わり、形がちがってしまう（赤）。}
\end{bkfigure}

もとの形の、どの長さも同じ倍率で大きくした図を\textbf{拡大図}、同じ倍率で小さくした図を\textbf{縮図}といいます。拡大図や縮図では、
\begin{itemize}[leftmargin=1.2\zw, itemsep=2pt, topsep=2pt, label={\color{\bkTT}●}]
\item 対応する（同じ場所どうしの）辺の長さは、すべて同じ倍率になります。
\item 対応する角の大きさは、もとと同じです。
\end{itemize}
横だけ3倍にした赤い家は、横の長さは3倍、たての長さは1倍で、倍率がそろっていません。だから拡大図ではありません。屋根の角を分度器ではかると、もとの家とちがっています。

\bkivdmidashi{縮尺：地図や設計図}
地図や建物の設計図は、本物の縮図です。本物をどれだけ縮めたかを表す比を、\textbf{縮尺}といいます。\bkidx{しゅくしゃく}{縮尺}縮尺 $1:\num{25000}$ の地図では、地図の上の\qty{1}{cm}が、本物の\num{25000}\,cmです。
\bkivdline{本物の長さ＝図の上の長さ×\num{25000}}
\bkivdline{図の上の長さ＝本物の長さ÷\num{25000}}
地図の上で\qty{4}{cm}なら、本物は、\qty{4}{cm}を\num{25000}倍して、\qty{100000}{cm}です。$\qty{1}{m}=\qty{100}{cm}$、$\qty{1}{km}=\qty{1000}{m}$ なので（2-10）、
\bkivdline{$\qty{100000}{cm}\div\qty{100}{cm/m}=\qty{1000}{m}$\qquad $\qty{1000}{m}\div\qty{1000}{m/km}=\qty{1}{km}$}
\qty{1}{km}です。かけ算のあとは、cmのままになっていることに注意しましょう。建物の設計図では、$1:100$ や $1:50$ の縮尺がよく使われます。

\bksubsection{2}{相似と相似比}

\leavevmode\label{bk:4-3:souji}2つの図形があって、一方を拡大または縮小すると、もう一方にぴったり重なるとき、2つの図形は\textbf{相似}（そうじ）であるといいます。\bkidx{そうじ}{相似}「形が同じで、大きさだけがちがう」ということです。やってみようの三角形ABCと三角形DEFは相似です。記号 $\sojisym$ を使って、
\bkivdline{$\triangle\mathrm{ABC}\sojisym\triangle\mathrm{DEF}$}
と書きます。$\sojisym$ は、英語で「にている」という意味の similar の頭文字Sを、横にねかせた形だといわれています。

\begin{bkfigure}
\begin{tikzpicture}[x=1mm, y=1mm]
  % ABC：A(0,0) B(24,0) C(0,18)（1 を 6 mm）。DEF は 1.5 倍：D(44,0) E(80,0) F(44,27)
  \fill[bkBlue!15] (0,0) -- (24,0) -- (0,18) -- cycle;
  \draw[bkBlue, line width=1.1pt] (0,0) -- (24,0) -- (0,18) -- cycle;
  \fill[bkOrange!15] (44,0) -- (80,0) -- (44,27) -- cycle;
  \draw[bkOrange, line width=1.1pt] (44,0) -- (80,0) -- (44,27) -- cycle;
  \node[font=\sffamily\footnotesize, anchor=north east] at (-0.4,-0.4) {A};
  \node[font=\sffamily\footnotesize, anchor=north west] at (24.4,-0.4) {B};
  \node[font=\sffamily\footnotesize, anchor=south] at (0,18.6) {C};
  \node[font=\sffamily\footnotesize, anchor=north east] at (43.6,-0.4) {D};
  \node[font=\sffamily\footnotesize, anchor=north west] at (80.4,-0.4) {E};
  \node[font=\sffamily\footnotesize, anchor=south] at (44,27.6) {F};
  \node[font=\sffamily\footnotesize, anchor=north] at (12,-1.6) {4};
  \node[font=\sffamily\footnotesize, anchor=south west] at (12.6,9.4) {5};
  \node[font=\sffamily\footnotesize, anchor=east] at (-1.4,9) {3};
  \node[font=\sffamily\footnotesize, anchor=north] at (62,-1.6) {6};
  \node[font=\sffamily\footnotesize, anchor=south west] at (62.6,14) {7.5};
  \node[font=\sffamily\footnotesize, anchor=east] at (42.6,13.5) {4.5};
  \bkivdright{(0,0)}{0}{90}
  \bkivdright{(44,0)}{0}{90}
  \node[font=\sffamily\small, align=left, anchor=west] at (90,14) {AとD、BとE、\\CとFが対応する\\（同じ場所の頂点）};
\end{tikzpicture}
\bkfigcaption{$\triangle\mathrm{ABC}\sojisym\triangle\mathrm{DEF}$。辺はどれも1.5倍。対応する角は等しい（たとえば$\angle\mathrm{A}$と$\angle\mathrm{D}$はどちらも直角）。}
\end{bkfigure}

相似の記号を書くときは、\textbf{対応する頂点を、同じ順番に書きます}。AとD、BとE、CとFが対応するなら、$\triangle\mathrm{ABC}\sojisym\triangle\mathrm{DEF}$ です。こう書いておけば、「1番めどうし・2番めどうし」で対応がすぐわかります。

\bkivdmidashi{相似比}
相似な図形の、対応する辺の長さの比を\textbf{相似比}といいます。\bkidx{そうじひ}{相似比}図の2つの三角形では、
\bkivdline{AB : DE＝\mbox{4 : 6＝2 : 3}\qquad \mbox{BC : EF＝5 : 7.5＝2 : 3}}
\bkivdline{CA : FD＝\mbox{3 : 4.5＝2 : 3}}
どの辺の比も $2:3$ なので、相似比は $2:3$ です。「大きい三角形は、小さい三角形の $\dfrac{3}{2}=1.5$ 倍」とも言えます。この倍率を $k$ で表すことがあります。ここでは $k=1.5$ です。

\bkivdmidashi{合同は、相似比 $1:1$ の相似}
合同（4-1）は、形も大きさも同じ図形でした。相似比が $1:1$ の相似が、合同です。合同の記号は $\equiv$ で、$\triangle\mathrm{ABC}\equiv\triangle\mathrm{DEF}$ のように、これも対応する順に書きます。相似の記号 $\sojisym$ は、アルファベットの S を横にたおした形です。なみの形の「～」とはちがい、両はしがくるりと丸まっています。合同の記号 $\equiv$ は、横線が3本です。（外国の本では、合同を $\cong$ と書くこともあります。）

\bkivdmidashi{1つの図形の中の比も、同じになる}
相似比は、2つの図形の「間」の比です。もう1つ、1つの図形の「中」の比も見てみましょう。三角形ABCで $\mathrm{AB}:\mathrm{CA}=4:3$、三角形DEFで $\mathrm{DE}:\mathrm{FD}=6:4.5=4:3$。同じです。形が同じなら、「横の辺はななめの辺の何倍か」のような、1つの図形の中の比も同じになります。この見方は、⑦と4-6で大切になります。

\bkivdmidashi{比の式から、長さを求める}
相似比が $2:3$ で、小さいほうの辺が\qty{8}{cm}なら、大きいほうの対応する辺 $x$\,cm は、$8:x=2:3$ です。比が等しいことは、比の値（2-5）が等しいことなので、
\bkivdline{$\dfrac{8}{x}=\dfrac{2}{3}$\qquad 両辺に $3x$ をかけて（2-14）\qquad $8\times3=2\times x$\qquad $x=12$}
「$a:b=c:d$ なら $a\times d=b\times c$（外側どうしの積＝内側どうしの積）」と、まとめて使えます。もっと簡単に、「大きいほうは $\dfrac{3}{2}$ 倍だから $8\times\dfrac{3}{2}=12$」と考えてもかまいません。

\bksubsection{3}{三角形の相似条件}

\leavevmode\label{bk:4-3:jouken}2つの三角形が相似かどうかを、3つの辺と3つの角を全部はかって確かめるのは大変です。合同のときは、「三角形が1つに決まる条件」が3つありました（4-1）。相似でも、「三角形の形が1つに決まる条件」を考えます。

\begin{bkcenter}
\begin{bktable}{colspec={Q[c,m,wd=8mm] Q[l,m,wd=50mm] Q[l,m,wd=56mm]}, rowsep=2pt, row{1}={font=\sffamily\bfseries\small}}
 & 合同条件（4-1） & 相似条件 \\
① & 3辺がそれぞれ等しい & 3辺の比がすべて等しい \\
② & 2辺とその間の角がそれぞれ等しい & 2辺の比とその間の角がそれぞれ等しい \\
③ & 1辺とその両はしの角がそれぞれ等しい & 2つの角がそれぞれ等しい \\
\end{bktable}
\bktabcaption{合同条件の「長さが等しい」を、「長さの比が等しい」にかえたものが相似条件。}
\end{bkcenter}

合同条件では、大きさまで決まります。相似条件では、形だけが決まり、大きさは決まりません。そのため、長さそのものではなく、長さの「比」が等しいことを条件にします。

\bkivdstep{1}{③「2つの角」だけで、なぜ形が決まるの？（図で考える）}
\noindent\begin{minipage}[c]{0.5\linewidth}\raggedright
1辺BCの長さを、2、3、4と変えて、その両はしにいつも40°と60°の角をかいてみます。できる三角形は、大きさはちがっても、すべて同じ形です。

辺の長さを自由に変えられるので、大きさは決まりません。でも、2つの角が決まると、3つめの角も $180^{\circ}-40^{\circ}-60^{\circ}=80^{\circ}$ に決まり（4-1）、形は1つに決まります。
\end{minipage}\hfill
\begin{minipage}[c]{0.46\linewidth}\centering
\begin{tikzpicture}[x=0.9mm, y=0.9mm]
  % B を原点に、BC＝20・30・40 mm。$\angle\mathrm{B}$＝40°、$\angle\mathrm{C}$＝60°。
  % 頂点 A までの長さ AB＝BC·sin60°/sin80°
  \foreach \L/\c in {40/bkPurple, 30/bkBlue, 20/bkOrange} {
    \pgfmathsetmacro\ab{\L*sin(60)/sin(80)}
    \draw[\c, line width=1pt] (0,0) -- (\L,0) -- (40:\ab) -- cycle;}
  \fill[bkInk] (0,0) circle (0.7);
  \draw[bkInk, line width=0.5pt] (6,0) arc (0:40:6);
  \node[font=\sffamily\footnotesize] at (20:9) {40°};
  \node[font=\sffamily\footnotesize, anchor=north] at (0,-1.4) {B};
  \foreach \L in {20,30,40} {\node[font=\sffamily\footnotesize, anchor=north] at (\L,-1.4) {C};}
  \node[font=\sffamily\footnotesize, anchor=south] at (38,8.4) {60°};
\end{tikzpicture}
\end{minipage}

\bkivdstep{2}{③「2つの角」だけで、なぜ形が決まるの？（数で考える）}
どの三角形でも、辺の長さをはかって、比を計算してみます（コンピュータで計算した値）。

\begin{bkcenter}
\begin{bktable}{colspec={Q[c,m,wd=16mm] Q[c,m,wd=18mm] Q[c,m,wd=18mm] Q[c,m,wd=18mm] Q[c,m,wd=34mm]}, rowsep=1.5pt, row{1}={font=\sffamily\small}}
BC & AB & CA & AB÷BC & CA÷BC \\
2 & 1.76 & 1.31 & 0.879 & 0.653 \\
3 & 2.64 & 1.96 & 0.879 & 0.653 \\
4 & 3.52 & 2.61 & 0.879 & 0.653 \\
\end{bktable}
\bktabcaption{BCを2倍にすると、ABもCAも2倍。辺の比はいつも同じ。}
\end{bkcenter}

BCが2から4に2倍になると、ABもCAも2倍です。どの大きさでも、AB÷BC は0.879、CA÷BC は0.653 で同じです。辺の比がすべて等しいので、相似です。

\bkivdmidashi{\bkhatenaref の答え}
\textbf{問い}\enspace 棒\qty{1}{m}の影が\qty{1.5}{m}、木の影が\qty{9}{m}。木の高さは？\par\vspace{2pt}
\noindent{\large\textbf{答え\enspace 木の高さは\qty{6}{m}}}\par\vspace{4pt}
\noindent\begin{minipage}[c]{0.56\linewidth}
\noindent{\sffamily\bfseries\color{bkOrange!75!black}わけ①\enspace 2つの三角形は相似}\par
太陽はとても遠いので、日光は平行にとどきます。木と棒は、どちらも地面に垂直に立っています。
\begin{itemize}[leftmargin=1.2\zw, itemsep=1pt, topsep=2pt, label={\color{\bkTT}●}]
\item 木と地面の角も、棒と地面の角も90°。
\item 日光と地面の角は、平行線の同位角なので等しい（4-1）。
\end{itemize}
2つの角がそれぞれ等しいので、2つの三角形は相似です（③）。
\end{minipage}\hfill
\begin{minipage}[c]{0.4\linewidth}\centering
\begin{tikzpicture}[x=0.9mm, y=0.9mm]
  % 1 m を 3.6 mm。木の三角形：高さ 21.6・影 32.4。棒：高さ 3.6・影 5.4（右に）
  \fill[bkGreen!15] (0,0) -- (32.4,0) -- (0,21.6) -- cycle;
  \draw[bkGreen!60!black, line width=1pt] (0,0) -- (0,21.6);
  \draw[bkInk, line width=0.8pt] (0,0) -- (32.4,0);
  \draw[bkOrange, line width=0.9pt] (0,21.6) -- (32.4,0);
  \bkivdright{(0,0)}{0}{90}
  \fill[bkGreen!15] (40,0) -- (45.4,0) -- (40,3.6) -- cycle;
  \draw[bkInk, line width=1.1pt] (40,0) -- (40,3.6);
  \draw[bkInk, line width=0.8pt] (40,0) -- (45.4,0);
  \draw[bkOrange, line width=0.9pt] (40,3.6) -- (45.4,0);
  \draw[bkInk, line width=0.5pt] (32.4,0) ++(180:5) arc (180:146.3:5);
  \draw[bkInk, line width=0.5pt] (45.4,0) ++(180:3) arc (180:146.3:3);
  \node[font=\sffamily\footnotesize, anchor=east] at (-1.4,10.8) {$x$\,m};
  \node[font=\sffamily\footnotesize, anchor=north] at (16.2,-1.4) {\qty{9}{m}};
  \node[font=\sffamily\footnotesize, anchor=north] at (42.7,-1.4) {\qty{1.5}{m}};
  \node[font=\sffamily\footnotesize, anchor=south] at (40,4.4) {\qty{1}{m}};
\end{tikzpicture}
\end{minipage}

\par\vspace{3pt}
\noindent{\sffamily\bfseries\color{bkOrange!75!black}わけ②\enspace 対応する辺の比で計算する}\par
対応する辺は、「高さ」どうしと「影」どうしです。影の比は $9:1.5=6:1$ なので、木の三角形は、棒の三角形を6倍に拡大した形です。だから、木の高さも棒の6倍で、$1\times6=6$（m）です。\par\vspace{2pt}
1つの三角形の中の比で考えることもできます。棒は「高さ÷影＝$1\div1.5$」なので、木も同じ比になり、木の高さは $9\times\dfrac{1}{1.5}=6$（m）です。どちらの考え方でも、答えは\qty{6}{m}です。

\bksubsection{4}{平行線と線分の比}

\leavevmode\label{bk:4-3:heikou}三角形ABCの辺AB、ACの上に点D、Eをとり、DEがBCに平行になるようにします。

\noindent\begin{minipage}[c]{0.44\linewidth}\centering
\begin{tikzpicture}[x=1mm, y=1mm]
  % A(16,36) B(0,0) C(44,0)。D・E は A から 2/3 の所
  \fill[bkBlue!10] (16,36) -- (0,0) -- (44,0) -- cycle;
  \fill[bkOrange!25] (16,36) -- ($(16,36)!0.6667!(0,0)$) -- ($(16,36)!0.6667!(44,0)$) -- cycle;
  \draw[bkInk, line width=0.9pt] (16,36) -- (0,0) -- (44,0) -- cycle;
  \draw[bkOrange, line width=1.2pt] ($(16,36)!0.6667!(0,0)$) -- ($(16,36)!0.6667!(44,0)$);
  \node[font=\sffamily\footnotesize, anchor=south] at (16,36.6) {A};
  \node[font=\sffamily\footnotesize, anchor=north east] at (-0.4,-0.4) {B};
  \node[font=\sffamily\footnotesize, anchor=north west] at (44.4,-0.4) {C};
  \node[font=\sffamily\footnotesize, anchor=east] at ($(16,36)!0.6667!(0,0)+(-0.8,0)$) {D};
  \node[font=\sffamily\footnotesize, anchor=west] at ($(16,36)!0.6667!(44,0)+(0.8,0)$) {E};
\end{tikzpicture}
\end{minipage}\hfill
\begin{minipage}[c]{0.52\linewidth}
DE∥BCなので、同位角が等しく（4-1）、$\angle\mathrm{ADE}$＝$\angle\mathrm{ABC}$、$\angle\mathrm{AED}$＝$\angle\mathrm{ACB}$ です。2つの角がそれぞれ等しいので、
\[ \triangle\mathrm{ADE}\sojisym\triangle\mathrm{ABC} \]
（相似条件③）。だから、対応する辺の比が等しくなります。
\end{minipage}
\bkivdline{\fbox{AD : AB＝\mbox{AE : AC＝DE : BC}}}
たとえば、AD＝4、DB＝2なら、AB＝$4+2=6$ です。だから、\mbox{AD : AB＝4 : 6＝2 : 3} です。DEは、BCの $\dfrac{2}{3}$ の長さになります。AEとECについても、\mbox{AE : EC＝AD : DB＝2 : 1} が成り立ちます。AEがACの $\dfrac{2}{3}$ なら、ECは残りの $\dfrac{1}{3}$ だからです。

\bkivdmidashi{X字形（砂時計の形）}
\noindent\begin{minipage}[c]{0.56\linewidth}
2本の直線が点Oで交わり、ABとDCが平行になっている形もよく出てきます。$\angle\mathrm{AOB}$＝$\angle\mathrm{COD}$（対頂角）、$\angle\mathrm{OAB}$＝$\angle\mathrm{OCD}$（平行線の錯角。4-1）なので、2つの角がそれぞれ等しく、
\[ \triangle\mathrm{OAB}\sojisym\triangle\mathrm{OCD} \]
です。三角形が、Oを中心にさかさまになって、大きさを変えた形です。ピンホールカメラやレンズは、この形で像を作ります（例題6、物理・工学ではこう使う）。
\end{minipage}\hfill
\begin{minipage}[c]{0.4\linewidth}\centering
\begin{tikzpicture}[x=1mm, y=1mm]
  % O(0,0)。A(-24,-12) B(-24,12)…左の大きい三角形。C(12,6) D(12,-6)…右の小さい三角形
  \fill[bkBlue!15] (0,0) -- (-24,12) -- (-24,-12) -- cycle;
  \fill[bkOrange!25] (0,0) -- (12,6) -- (12,-6) -- cycle;
  \draw[bkInk, line width=0.9pt] (-24,12) -- (12,-6);
  \draw[bkInk, line width=0.9pt] (-24,-12) -- (12,6);
  \draw[bkBlue, line width=1.2pt] (-24,12) -- (-24,-12);
  \draw[bkOrange, line width=1.2pt] (12,6) -- (12,-6);
  \fill[bkInk] (0,0) circle (0.8);
  \node[font=\sffamily\footnotesize, anchor=north] at (0,-1.4) {O};
  \node[font=\sffamily\footnotesize, anchor=east] at (-24.6,-12) {A};
  \node[font=\sffamily\footnotesize, anchor=east] at (-24.6,12) {B};
  \node[font=\sffamily\footnotesize, anchor=west] at (12.6,6) {C};
  \node[font=\sffamily\footnotesize, anchor=west] at (12.6,-6) {D};
\end{tikzpicture}
\end{minipage}

\begin{column}{5円玉の穴で、月の大きさをはかる}
\noindent\begin{minipage}[c]{0.6\linewidth}\raggedright
5円玉の穴の直径は、約\qty{5}{mm}です。満月の夜、5円玉を目の前に持ち、ちょうど穴いっぱいに月が見える所をさがします。目から5円玉までが約\qty{55}{cm}のとき、月がぴったり穴におさまりました。

目と穴のふち、目と月のふちを結ぶと、相似な2つの三角形ができます（④の平行線と線分の比）。月までの距離は約\qty{384000}{km}なので、相似比は $\qty{55}{cm}:\qty{384000}{km}$ です。単位をmにそろえると、月の直径は
\[ \qty{0.005}{m}\times\frac{\qty{384000000}{m}}{\qty{0.55}{m}}\approx\qty{3490000}{m} \]
で、約\qty{3490}{km}。本当の月の直径は約\qty{3474}{km}なので、5円玉1まいで、ほぼ正しく求められます。
\end{minipage}\hfill
\begin{minipage}[c]{0.36\linewidth}\centering
\begin{tikzpicture}[x=0.72mm, y=0.72mm]
  % 目 (0,0)、5円玉 (12,0)、月 (44,0)（縮尺は合っていない）
  \fill[bkInk] (0,0) circle (1);
  \node[font=\sffamily\footnotesize, anchor=north] at (0,-1.8) {目};
  \draw[bkYellow!60!black, line width=1.6pt] (12,-7.5) -- (12,-2.73) (12,2.73) -- (12,7.5);
  \node[font=\sffamily\footnotesize, anchor=north] at (12,-8.5) {5円玉};
  \fill[bkYellow!40] (44,0) circle (10);
  \draw[bkYellow!60!black, line width=0.8pt] (44,0) circle (10);
  \node[font=\sffamily\footnotesize, anchor=north] at (44,-11.5) {月};
  \draw[bkOrange, line width=0.6pt] (0,0) -- (44,10) (0,0) -- (44,-10);
\end{tikzpicture}
\end{minipage}
\end{column}

\bksubsection{5}{面積の比は $k^2$、体積の比は $k^3$}

\leavevmode\label{bk:4-3:menseki}相似比が $1:2$ のとき、長さは2倍です。では、面積は何倍になるでしょう。「2倍」と答えたくなりますが、方眼で確かめてみましょう。

\begin{bkfigure}
\begin{tikzpicture}[x=1mm, y=1mm]
  % 正方形 1辺 1・2・3（1マス 7 mm）
  \draw[bkBlue, line width=0.9pt, fill=bkBlue!20] (0,0) rectangle (7,7);
  \begin{scope}[shift={(18,0)}]
    \fill[bkBlue!20] (0,0) rectangle (14,14);
    \draw[bkBlue!60, line width=0.4pt] (0,0) grid[step=7] (14,14);
    \draw[bkBlue, line width=0.9pt] (0,0) rectangle (14,14);
  \end{scope}
  \begin{scope}[shift={(43,0)}]
    \fill[bkBlue!20] (0,0) rectangle (21,21);
    \draw[bkBlue!60, line width=0.4pt] (0,0) grid[step=7] (21,21);
    \draw[bkBlue, line width=0.9pt] (0,0) rectangle (21,21);
  \end{scope}
  \node[font=\sffamily\footnotesize, anchor=north] at (3.5,-1.6) {1こ};
  \node[font=\sffamily\footnotesize, anchor=north] at (25,-1.6) {$2\times2=4$こ};
  \node[font=\sffamily\footnotesize, anchor=north] at (53.5,-1.6) {$3\times3=9$こ};
  % 立方体 1辺 1・2（1マス 6 mm）。斜め投影
  \begin{scope}[shift={(78,0)}]
    \foreach \x/\y/\z in {0/0/0} {
      \fill[bkOrange!30] (0,0) rectangle (6,6);
      \fill[bkOrange!50] (6,0) -- (9,3) -- (9,9) -- (6,6) -- cycle;
      \fill[bkOrange!15] (0,6) -- (6,6) -- (9,9) -- (3,9) -- cycle;
      \draw[bkOrange!80!black, line width=0.6pt] (0,0) rectangle (6,6) (6,0) -- (9,3) -- (9,9) -- (3,9) -- (0,6) (6,6) -- (9,9);}
    \node[font=\sffamily\footnotesize, anchor=north] at (4.5,-1.6) {1こ};
  \end{scope}
  \begin{scope}[shift={(96,0)}]
    \fill[bkOrange!30] (0,0) rectangle (12,12);
    \fill[bkOrange!50] (12,0) -- (18,6) -- (18,18) -- (12,12) -- cycle;
    \fill[bkOrange!15] (0,12) -- (12,12) -- (18,18) -- (6,18) -- cycle;
    \draw[bkOrange!70!black, line width=0.35pt] (6,0) -- (6,12) (0,6) -- (12,6) (12,6) -- (18,12) (15,3) -- (15,15) (6,12) -- (12,18) (3,15) -- (15,15);
    \draw[bkOrange!80!black, line width=0.6pt] (0,0) rectangle (12,12) (12,0) -- (18,6) -- (18,18) -- (6,18) -- (0,12) (12,12) -- (18,18);
    \node[font=\sffamily\footnotesize, anchor=north] at (9,-1.6) {$2\times2\times2=8$こ};
  \end{scope}
\end{tikzpicture}
\bkfigcaption{左：正方形の1辺を2倍、3倍にすると、面積は4倍、9倍。右：立方体の1辺を2倍にすると、体積は8倍。}
\end{bkfigure}

1辺を2倍にした正方形には、もとの正方形が、たてに2こ、横に2こで、$2\times2=4$ こ入ります。たても横も2倍になるので、2が2回かかるのです。1辺を3倍にすると、$3\times3=9$ 倍です。

\bkivdstep{1}{どんな形でも、面積は $k^2$ 倍}
三角形なら、面積は「底辺×高さ÷2」です。相似比が $1:k$ なら、底辺も高さも $k$ 倍なので、面積は $k\times k=k^2$ 倍です。曲がった形でも、細かい方眼のマスに分けて考えます（4-2の⑤と同じ「細かく分ける」考え方）。$k$ 倍に拡大すると、マスの1つ1つが $k^2$ 倍になるので、全体も $k^2$ 倍です。

\bkivdstep{2}{体積は $k^3$ 倍}
立方体の1辺を2倍にすると、たて・横・高さのすべてが2倍なので、小さい立方体が $2\times2\times2=8$ こ入ります。2が3回かかります。どんな立体でも、相似比が $1:k$ なら、体積は $k^3$ 倍です。

\begin{bkcenter}
\begin{bktable}{colspec={Q[l,m,wd=36mm] *{5}{Q[c,m,wd=14mm]}}, rowsep=2.5pt}
長さの倍率 $k$ & 1 & 2 & 3 & 10 & $\dfrac{1}{2}$ \\
面積の倍率 $k^2$ & 1 & 4 & 9 & 100 & $\dfrac{1}{4}$ \\
体積の倍率 $k^3$ & 1 & 8 & 27 & 1000 & $\dfrac{1}{8}$ \\
\end{bktable}
\bktabcaption{相似比が $1:k$ のとき、面積の比は $1:k^2$、体積の比は $1:k^3$。}
\end{bkcenter}

\begin{taisetsu}
相似比が $1:k$ のとき\\
長さの比 $1:k$\qquad 面積の比 $1:k^2$\qquad 体積の比 $1:k^3$
\end{taisetsu}

\begin{column}{コピー機の「141\%」と、紙の大きさ}
\noindent\begin{minipage}[c]{0.64\linewidth}\raggedright
コピー機でA4の紙をA3の紙に拡大するときは、倍率を141\%（1.41倍）にします。なぜ2倍ではないのでしょう。

A3の紙は、A4の紙2まいをつなげた大きさで、面積がちょうど2倍です。A4とA3は相似な長方形になるように決めてあります。面積の倍率 $k^2$ が2なので、長さの倍率は $k=\sqrt{2}=1.414\cdots$（2-8）です。だから、長さを約1.41倍にすると、ぴったりA3の大きさになります。反対に、A3をA4に縮小するときは、$\dfrac{1}{\sqrt{2}}\approx0.71$ で、71\%です（$\approx$ は「ほぼ等しい」という意味。5-8でくわしく）。

本当の大きさは、A4が \qty{210}{mm}×\qty{297}{mm}、A3が \qty{297}{mm}×\qty{420}{mm}。$297\div210\approx1.414$、$420\div297\approx1.414$ で、確かに $\sqrt{2}$ 倍です。
\end{minipage}\hfill
\begin{minipage}[c]{0.32\linewidth}\centering
\begin{tikzpicture}[x=1mm, y=1mm]
  % A3（たて 42・横 29.7 を 0.8 倍）を、横線で A4 2まいに分ける
  \fill[bkBlue!15] (0,0) rectangle (23.76,33.6);
  \draw[bkBlue!80!black, line width=0.9pt] (0,0) rectangle (23.76,33.6);
  \draw[bkBlue!80!black, line width=0.7pt, dash pattern=on 2pt off 1pt] (0,16.8) -- (23.76,16.8);
  \node[font=\sffamily\footnotesize] at (11.88,8.4) {A4};
  \node[font=\sffamily\footnotesize] at (11.88,25.2) {A4};
  \node[font=\sffamily\footnotesize\bfseries, anchor=south] at (11.88,34.4) {A3};
\end{tikzpicture}
\end{minipage}
\end{column}

\par\Needspace{10\baselineskip}\bkivdmidashi{大きい動物の足が太いわけ}
\noindent\begin{minipage}[c]{0.62\linewidth}
シカをそのまま3倍の大きさにした動物を考えます。体重は体積に比例するので、$3^3=27$ 倍になります。体重をささえる足の骨の強さは、骨の断面の面積に比例するので、$3^2=9$ 倍にしかなりません。骨の断面1\,\unit{cm^2}あたりにかかる重さは、$27\div9=3$ 倍になってしまいます。だから、ゾウのような大きな動物は、体の大きさのわりに、ずっと太い足をしています。

このように「大きさを変えると、面積で決まる量と体積で決まる量の割合が変わる」ことを調べるのが、物理や工学の\textbf{相似則}です。
\end{minipage}\hfill
\begin{minipage}[c]{0.34\linewidth}\centering
\begin{tikzpicture}[x=1mm, y=1mm]
  % 足の断面（円）と体（立方体）の大きさのくらべ
  \draw[bkInk!60, line width=0.6pt, fill=bkStem!30] (0,0) circle (2);
  \draw[bkInk!60, line width=0.6pt, fill=bkStem!30] (16,0) circle (6);
  \node[font=\sffamily\footnotesize, anchor=north] at (0,-7) {断面 1};
  \node[font=\sffamily\footnotesize, anchor=north] at (16,-7) {断面 9};
  \draw[bkInk!60, line width=0.6pt, fill=bkOrange!25] (-3,14) rectangle (3,20);
  \draw[bkInk!60, line width=0.6pt, fill=bkOrange!25] (7,14) rectangle (25,32);
  \node[font=\sffamily\footnotesize, anchor=north] at (0,13) {体重 1};
  \node[font=\sffamily\footnotesize] at (16,23) {体重 27};
\end{tikzpicture}
\end{minipage}

\begin{motto}[3]
\textbf{光の明るさと、距離の2乗}\par
プロジェクターでスクリーンに映すとき、スクリーンまでの距離を2倍にすると、映る画面は相似なまま、たても横も2倍になります。面積は $2^2=4$ 倍です。同じ量の光が4倍の面積に広がるので、画面の明るさは $\dfrac{1}{4}$ になります。電球や星の光も同じで、「光の明るさは、距離の2乗に反比例する」のは、この相似の $k^2$ が理由です。物理の音の強さや、重力・電気の力が距離の2乗に反比例することにも、同じ考え方が出てきます。
\end{motto}

\bksubsection{6}{すいを底面に平行に切ると}

\leavevmode\label{bk:4-3:sui}アイスクリームのコーンや、とんがり帽子のような形を\textbf{円すい}、ピラミッドのような、底面が四角形の形を\textbf{四角すい}といいます。とがった点を頂点といいます。円すいを、底面に平行な平面で切ってみましょう。切り口（断面）は円です。その半径は、どこで切ると、どれだけになるでしょう。

\noindent\begin{minipage}[c]{0.44\linewidth}\centering
\begin{tikzpicture}[x=1mm, y=1mm]
  % 円すい：頂点 (0,36)、底面の中心 (0,0)、半径 12（高さ 12 cm・半径 4 cm を 3 倍）
  \fill[bkBlue!12] (0,36) -- (-12,0) arc (180:360:12 and 3) -- cycle;
  \draw[bkBlue, line width=0.9pt] (-12,0) arc (180:360:12 and 3);
  \draw[bkBlue, line width=0.6pt, dash pattern=on 1.5pt off 1.5pt] (12,0) arc (0:180:12 and 3);
  \draw[bkBlue, line width=0.9pt] (-12,0) -- (0,36) -- (12,0);
  \fill[bkOrange!35] (0,18) ellipse (6 and 1.5);
  \draw[bkOrange!80!black, line width=0.9pt] (0,18) ellipse (6 and 1.5);
  \draw[bkGray, line width=0.6pt, dash pattern=on 2pt off 1pt] (0,36) -- (0,0);
  \draw[bkInk, line width=0.7pt] (0,18) -- (6,18);
  \draw[bkInk, line width=0.7pt] (0,0) -- (12,0);
  \fill[bkInk] (0,36) circle (0.8);
  \node[font=\sffamily\footnotesize, anchor=south] at (0,36.8) {頂点};
  \node[font=\sffamily\footnotesize, anchor=west] at (13.5,0) {半径 $r$};
  \node[font=\sffamily\footnotesize, anchor=west] at (8.5,18) {断面の半径};
  \draw[bkInk, line width=0.5pt, {Stealth[length=2.5pt]}-{Stealth[length=2.5pt]}] (-15.5,36) -- (-15.5,18);
  \node[font=\sffamily\footnotesize, anchor=east] at (-16.5,27) {$y$};
  \draw[bkInk, line width=0.5pt, {Stealth[length=2.5pt]}-{Stealth[length=2.5pt]}] (-21.5,36) -- (-21.5,0);
  \node[font=\sffamily\footnotesize, anchor=east] at (-22.5,13) {$h$};
\end{tikzpicture}
\end{minipage}\hfill
\begin{minipage}[c]{0.52\linewidth}
円すいを、頂点を通るたての平面で切って、横から見ます。頂点から $y$ の所で、底面に平行に切った線は、三角形の底辺に平行です。

だから、④の平行線と線分の比が使えます。頂点からの長さの比 $y:h$ と、半径の比が等しくなります。高さ $h$ で半径 $r$ の円すいなら、
\[ \text{断面の半径}=r\times\frac{y}{h}=\frac{r}{h}\,y \]
\end{minipage}

\par\vspace{3pt}
$\dfrac{r}{h}$ は、$y$ が変わっても変わらない数（定数）です。だから、\textbf{断面の半径は、頂点からの長さ $y$ に比例します}（2-15）。頂点から切りとった小さい円すいは、もとの円すいと相似で、相似比は $y:h$ です。高さ\qty{12}{cm}、底面の半径\qty{4}{cm}の円すいで計算すると、次のようになります。

\begin{bkcenter}
\begin{bktable}{colspec={Q[l,m,wd=40mm] *{4}{Q[c,m,wd=16mm]}}, rowsep=1.5pt}
頂点からの長さ $y$（cm） & 3 & 6 & 9 & 12 \\
断面の半径（cm） & 1 & 2 & 3 & 4 \\
断面の面積（\unit{cm^2}） & $\pi$ & $4\pi$ & $9\pi$ & $16\pi$ \\
\end{bktable}
\bktabcaption{半径は $y$ に比例（$\dfrac{4}{12}y=\dfrac{1}{3}y$）。面積は $y^2$ に比例する（⑤の $k^2$）。}
\end{bkcenter}

四角すいでも同じです。底面が1辺 $a$ の正方形で、高さが $h$ の四角すいを、頂点から $y$ の所で底面に平行に切ると、断面は1辺 $\dfrac{a}{h}\,y$ の正方形です。第5巻の5-10では、この断面をうすい板にして積み重ね、すいの体積を積分で求めます。

\bksubsection{7}{直角三角形の辺の比は、角だけで決まる}

\leavevmode\label{bk:4-3:chokkaku}直角三角形では、1つの角は必ず90°です。だから、もう1つの角 $\theta$ が同じなら、2つの角がそれぞれ等しいので、相似です（③）。相似なら、1つの三角形の中の辺の比も同じでした（②）。
\bkivdline{\textbf{直角三角形の辺の比は、直角でない1つの角 $\theta$ だけで決まる}}

\noindent\begin{minipage}[c]{0.48\linewidth}\centering
\begin{tikzpicture}[x=1mm, y=1mm]
  % θ＝30°の直角三角形。斜辺 20・36・52 mm
  \foreach \c/\hyp in {bkPurple/52, bkBlue/36, bkOrange/20} {
    \pgfmathsetmacro\bx{\hyp*cos(30)}\pgfmathsetmacro\hy{\hyp*sin(30)}
    \draw[\c, line width=1pt] (0,0) -- (\bx,0) -- (\bx,\hy) -- cycle;
    \bkivdright{(\bx,0)}{180}{90}}
  \draw[bkInk, line width=0.5pt] (8,0) arc (0:30:8);
  \node[font=\sffamily\footnotesize] at (9:12.5) {30°};
  \node[font=\sffamily\footnotesize, anchor=north west] at (-1,-1.4) {斜辺の長さ：2、3.6、5.2};
\end{tikzpicture}
\end{minipage}\hfill
\begin{minipage}[c]{0.48\linewidth}
図は、30°の角をもつ直角三角形を、3つの大きさでかいたものです。直角に向かい合う、いちばん長い辺を\textbf{斜辺}といいます。どの三角形でも、
\bkivdline{たての辺÷斜辺＝0.5}
です。斜辺が2倍になれば、たての辺も2倍になるからです。
\end{minipage}

\bkivdmidashi{30°の三角形で、たての辺が斜辺のちょうど半分になるわけ}
\noindent\begin{minipage}[c]{0.36\linewidth}\centering
\begin{tikzpicture}[x=1mm, y=1mm]
  % 正三角形（1辺 32）を半分に
  \fill[bkBlue!15] (0,0) -- (27.71,16) -- (0,32) -- cycle;
  \fill[bkOrange!30] (0,0) -- (27.71,16) -- (0,16) -- cycle;
  \draw[bkInk, line width=0.9pt] (0,0) -- (27.71,16) -- (0,32) -- cycle;
  \draw[bkOrange!80!black, line width=1pt] (0,16) -- (27.71,16);
  \bkivdright{(0,16)}{0}{-90}
  \draw[bkInk, line width=0.5pt] (27.71,16) ++(180:6) arc (180:210:6);
  \node[font=\sffamily\footnotesize] at ($(27.71,16)+(195:9)$) {30°};
  \node[font=\sffamily\footnotesize, anchor=east] at (-1.2,8) {16};
  \node[font=\sffamily\footnotesize, anchor=east] at (-1.2,24) {16};
  \node[font=\sffamily\footnotesize, anchor=north west] at (15,6.2) {32};
\end{tikzpicture}
\end{minipage}\hfill
\begin{minipage}[c]{0.6\linewidth}
1辺32の正三角形は、3つの角がどれも60°です。これを、まん中の線で半分に折ると、2つの直角三角形ができます（4-1の線対称）。

1つの直角三角形（オレンジ）の角は、90°と60°と、$60^{\circ}\div2=30^{\circ}$ です。たての辺は、正三角形の1辺32の半分の16。斜辺は32のままです。だから、30°の直角三角形では、たての辺はいつも斜辺の半分です。
\end{minipage}

\bkivdmidashi{坂の傾きも、角だけで決まる}
3-5では、坂の急さを、傾き（たての変化量÷横の変化量）で表しました。坂の角が同じなら、直角三角形は相似なので、坂の長さに関係なく、傾き「たて÷横」も同じになります。たとえば、角が20°の坂では、横に\qty{10}{m}進むと約\qty{3.64}{m}上がり、横に\qty{50}{m}進むと約\qty{18.2}{m}上がります。どちらも、たて÷横＝0.364 です。

このように、直角三角形の「たて÷斜辺」「横÷斜辺」「たて÷横」は、角 $\theta$ を決めると、1つに決まる数です。4-6では、この3つの比に $\sin\theta$・$\cos\theta$・$\tan\theta$ という名前をつけて、坂や力の計算に使います。

\begin{taisetsu}
相似：一方を拡大・縮小すると、もう一方に重なる。対応する辺の比はすべて等しく、対応する角は等しい。\par\vspace{2pt}
三角形の相似条件：① 3辺の比がすべて等しい　② 2辺の比とその間の角が等しい　③ 2つの角が等しい\par\vspace{2pt}
DE∥BC なら　\mbox{AD : AB＝AE : AC＝DE : BC}\par\vspace{2pt}
相似比 $1:k$ なら、面積の比 $1:k^2$、体積の比 $1:k^3$\par\vspace{2pt}
すいを底面に平行に切った断面の半径（1辺）は、頂点からの長さに比例する。\par\vspace{2pt}
直角三角形の辺の比は、直角でない1つの角だけで決まる。
\end{taisetsu}

\begin{kotoba}
\small
\bkword{拡大図・縮図}{どの長さも同じ倍率で大きく（小さく）した図。角の大きさはもとと同じ。}
\bkword{縮尺}{本物の長さを、図の上でどれだけ縮めたかを表す比。$1:\num{25000}$ など。}
\bkword{相似・相似比}{形が同じで大きさがちがう図形どうしが相似。対応する辺の長さの比が相似比。}
\bkword{相似条件}{三角形が相似になるための条件。3辺の比、2辺の比とその間の角、2つの角。}
\bkword{相似則}{大きさを変えたとき、長さ・面積・体積で決まる量がそれぞれ何倍になるかを使って、本物を調べる考え方。}
\bkword{斜辺}{直角三角形で、直角に向かい合う、いちばん長い辺。}
\end{kotoba}

\bkkigocards{\bkkigocard{$\sojisym$}{そうじ}{相似。$\triangle\mathrm{ABC}\sojisym\triangle\mathrm{DEF}$ のように、対応する頂点の順に書く}{similar}{4-3}}
{\bkkigocard{$\equiv$}{ごうどう}{合同。形も大きさも同じ。相似比 $1:1$ の相似}{congruent}{4-3}}

%% ---------------- 例題 ----------------
\bksection{reidai}{例題でたしかめよう}

\begin{reidai}[縮尺から長さを求める]
縮尺 $1:\num{25000}$ の地図があります。（1）\nobreak 地図の上で\qty{6}{cm}の道のりは、本物では何kmですか。（2）\nobreak 本物で\qty{2}{km}の道のりは、地図の上で何cmですか。\par
\kangaekata 本物は、地図の\num{25000}倍です。単位はcmのまま計算して、最後にkmに直します。\par
\shiki （1）$\qty{6}{cm}\times\num{25000}=\qty{150000}{cm}$\par
\hspace*{4.3\zw}$\qty{150000}{cm}\div\qty{100}{cm/m}=\qty{1500}{m}$\par
\hspace*{4.3\zw}$\qty{1500}{m}\div\qty{1000}{m/km}=\qty{1.5}{km}$\par
\hspace*{2.6\zw}（2）$\qty{2}{km}=\qty{2000}{m}=\qty{200000}{cm}$\par
\hspace*{4.3\zw}$\qty{200000}{cm}\div\num{25000}=\qty{8}{cm}$\par
\kotae （1）\nobreak \qty{1.5}{km}\quad（2）\nobreak \qty{8}{cm}
\end{reidai}

\begin{jibun}
縮尺 $1:50$ の家の設計図があります。（1）\nobreak 図の上で\qty{9}{cm}の長さは、本物では何mですか。（2）\nobreak 本物で\qty{3}{m}の長さは、図の上で何cmですか。
\bkjans{（1）\nobreak $9\times50=\qty{450}{cm}$ で\qty{4.5}{m}。（2）\nobreak \qty{3}{m}＝\qty{300}{cm}。$300\div50=\qty{6}{cm}$。}
\end{jibun}
\bkivdpoint{縮尺の計算は、① 長さの単位をcmにそろえる、② かける・わる、③ 最後にmやkmに直す、の順にします。}

\begin{reidai}[相似比から辺の長さを求める]
$\triangle\mathrm{ABC}\sojisym\triangle\mathrm{DEF}$ で、AB＝\qty{6}{cm}、BC＝\qty{8}{cm}、CA＝\qty{5}{cm}、DE＝\qty{9}{cm}です。相似比と、EF、FDの長さを求めましょう。\par
\kangaekata 記号の順番から、ABとDE、BCとEF、CAとFDが対応します。長さがわかっている対応する辺 AB、DE から、相似比を求めます。\par
\begin{bkcenter}
\begin{tikzpicture}[x=1mm, y=1mm]
  % B(0,0) C(8,0) A(4.6875,3.745) を 3.6 倍。DEF はその 1.5 倍（E が左はし）
  \fill[bkBlue!15] (0,0) -- (28.8,0) -- (16.88,13.48) -- cycle;
  \draw[bkBlue, line width=1pt] (0,0) -- (28.8,0) -- (16.88,13.48) -- cycle;
  \node[font=\sffamily\footnotesize, anchor=north east] at (-0.3,-0.3) {B};
  \node[font=\sffamily\footnotesize, anchor=north west] at (29.1,-0.3) {C};
  \node[font=\sffamily\footnotesize, anchor=south] at (16.88,14.1) {A};
  \node[font=\sffamily\footnotesize, anchor=south east] at (8.2,7.2) {6};
  \node[font=\sffamily\footnotesize, anchor=south west] at (23.2,7.2) {5};
  \node[font=\sffamily\footnotesize, anchor=north] at (14.4,-1.4) {8};
  \begin{scope}[shift={(48,0)}]
    \fill[bkOrange!20] (0,0) -- (43.2,0) -- (25.31,20.22) -- cycle;
    \draw[bkOrange, line width=1pt] (0,0) -- (43.2,0) -- (25.31,20.22) -- cycle;
    \node[font=\sffamily\footnotesize, anchor=north east] at (-0.3,-0.3) {E};
    \node[font=\sffamily\footnotesize, anchor=north west] at (43.5,-0.3) {F};
    \node[font=\sffamily\footnotesize, anchor=south] at (25.31,20.8) {D};
    \node[font=\sffamily\footnotesize, anchor=south east] at (12.2,10.8) {9};
    \node[font=\sffamily\footnotesize, anchor=south west] at (34.9,10.8) {？};
    \node[font=\sffamily\footnotesize, anchor=north] at (21.6,-1.4) {？};
  \end{scope}
\end{tikzpicture}
\end{bkcenter}

\shiki 相似比は $\mathrm{AB}:\mathrm{DE}=6:9=2:3$（3でわる）\par
\hspace*{2.6\zw}$\triangle\mathrm{DEF}$は$\triangle\mathrm{ABC}$の $\dfrac{3}{2}$ 倍。$\mathrm{EF}=8\times\dfrac{3}{2}=12$、$\mathrm{FD}=5\times\dfrac{3}{2}=7.5$\par
\hspace*{2.6\zw}確かめ：$8:x=2:3$ なら $2x=24$、$x=12$\par
\kotae 相似比 $2:3$、EF＝\qty{12}{cm}、FD＝\qty{7.5}{cm}
\end{reidai}

\begin{jibun}
$\triangle\mathrm{ABC}\sojisym\triangle\mathrm{DEF}$ で、AB＝\qty{4}{cm}、BC＝\qty{6}{cm}、CA＝\qty{5}{cm}、DE＝\qty{10}{cm}です。相似比と、EF、FDの長さを求めましょう。
\bkjans{相似比 $4:10=2:5$。$\mathrm{EF}=6\times\dfrac{5}{2}=\qty{15}{cm}$、$\mathrm{FD}=5\times\dfrac{5}{2}=\qty{12.5}{cm}$。}
\end{jibun}
\bkivdpoint{相似比は、長さがどちらもわかっている「対応する辺」の組から作ります。どれとどれが対応するかは、記号の頂点の順番がヒントです。}

\begin{reidai}[相似条件を見つける]
次の2つの三角形は相似ですか。相似なら、使った相似条件と相似比を答えましょう。\par
（1）\nobreak 3辺が 4、6、8 の三角形と、3辺が 6、9、12 の三角形\par
（2）\nobreak 2辺が 3、4 でその間の角が60°の三角形と、2辺が 6、8 でその間の角が60°の三角形\par
（3）\nobreak 2つの角が50°、70°の三角形と、2つの角が60°、70°の三角形\par
\kangaekata 短い辺どうし、長い辺どうしを組にして、比を調べます。角だけのときは、3つめの角も計算します。\par
\shiki （1）$4:6=2:3$、$6:9=2:3$、$8:12=2:3$。3辺の比がすべて等しい\par
\hspace*{2.6\zw}（2）$3:6=1:2$、$4:8=1:2$ で、その間の角はどちらも60°\par
\hspace*{2.6\zw}（3）1つめの3つめの角は $180^{\circ}-50^{\circ}-70^{\circ}=60^{\circ}$\par
\hspace*{4.3\zw}2つめの3つめの角は $180^{\circ}-60^{\circ}-70^{\circ}=50^{\circ}$\par
\hspace*{4.3\zw}どちらも50°、60°、70°の三角形。2つの角が等しい\par
\kotae （1）\nobreak 相似。3辺の比、相似比 $2:3$\quad（2）\nobreak 相似。2辺の比とその間の角、相似比 $1:2$\quad（3）\nobreak 相似。2つの角（相似比は、角だけではわからない）
\end{reidai}

\begin{jibun}
次の2つの三角形は相似ですか。相似なら、相似条件と相似比を答えましょう。（1）\nobreak 3辺が 3、5、6 と、6、10、12\quad（2）\nobreak 2辺が 5、6 でその間の角が40°と、2辺が 15、18 でその間の角が40°\quad（3）\nobreak 2つの角が45°、80°と、55°、80°
\bkjans{（1）\nobreak 相似。3辺の比、$1:2$。（2）\nobreak 相似。2辺の比とその間の角、$1:3$。（3）\nobreak 3つめの角は55°と45°。どちらも45°、55°、80°なので相似（2つの角）。}
\end{jibun}
\bkivdpoint{角で調べるときは、3つめの角を180°からひいて出してから比べます。辺の比は、短い辺どうし・まん中どうし・長い辺どうしを組にします。}

\begin{reidai}[影の長さから高さを求める]
長さ\qty{1.2}{m}の棒をまっすぐ立てると、影の長さは\qty{1.8}{m}でした。同じ時刻に、ビルの影の長さは\qty{15}{m}でした。ビルの高さを求めましょう。\par
\kangaekata 日光は平行で、棒もビルも地面に垂直なので、2つの三角形は相似です（2つの角）。「高さ÷影」が同じになります。\par
\vspace{2pt}\begin{bkcenter}
\begin{tikzpicture}[x=1mm, y=1mm]
  % 1 m を 2.2 mm。ビル 10 m（22 mm）・影 15 m（33 mm）。棒 1.2 m・影 1.8 m は 5 倍に拡大してかく
  \draw[bkGray, line width=0.8pt] (-4,0) -- (100,0);
  \fill[bkGray!35] (0,0) rectangle (9,22);
  \draw[bkInk!70, line width=0.7pt] (0,0) rectangle (9,22);
  \foreach \yy in {4,9,14,19} {\fill[white] (2,\yy) rectangle (4,\yy+2.4); \fill[white] (5.5,\yy) rectangle (7.5,\yy+2.4);}
  \fill[bkInk!50] (9,-0.8) rectangle (42,0);
  \draw[bkOrange, line width=0.7pt, dash pattern=on 2pt off 1.2pt] (9,22) -- (42,0);
  \node[font=\sffamily\footnotesize, anchor=north] at (25.5,-1.4) {影 \qty{15}{m}};
  \node[font=\sffamily\footnotesize, anchor=east] at (-1,11) {？};
  \draw[bkInk, line width=1.4pt] (70,0) -- (70,13.2);
  \fill[bkInk!50] (70,-0.8) rectangle (89.8,0);
  \draw[bkOrange, line width=0.7pt, dash pattern=on 2pt off 1.2pt] (70,13.2) -- (89.8,0);
  \node[font=\sffamily\footnotesize, anchor=north] at (79.9,-1.4) {影 \qty{1.8}{m}};
  \node[font=\sffamily\footnotesize, anchor=east] at (69,6.6) {\qty{1.2}{m}};
  \node[font=\sffamily\footnotesize, text=bkGray, anchor=west] at (90,9) {（拡大図）};
\end{tikzpicture}
\end{bkcenter}

\shiki 棒の高さ÷影＝$1.2\div1.8=\dfrac{12}{18}=\dfrac{2}{3}$\par
\hspace*{2.6\zw}ビルの高さ＝$15\times\dfrac{2}{3}=10$\par
\hspace*{2.6\zw}確かめ：影の比 $15:1.8$ と高さの比 $10:1.2$ は、どちらも $\dfrac{25}{3}$ 倍\par
\kotae \qty{10}{m}
\end{reidai}

\begin{jibun}
長さ\qty{1}{m}の棒の影が\qty{0.8}{m}のとき、同じ時刻に影の長さが\qty{6.4}{m}の木の高さを求めましょう。
\bkjans{高さ÷影＝$1\div0.8=1.25$。$6.4\times1.25=\qty{8}{m}$。}
\end{jibun}
\bkivdpoint{影の問題では、「同じ時刻」なら日光のかたむきが同じです。だから「高さ÷影」が同じ数になります。}

\begin{reidai}[平行線と線分の比]
\noindent\begin{minipage}[t]{0.62\linewidth}
三角形ABCで、辺AB上の点D、辺AC上の点Eを、DE∥BCとなるようにとります。AD＝4、DB＝2、AE＝6、BC＝9のとき、ECとDEの長さを求めましょう。\par
\kangaekata $\triangle\mathrm{ADE}\sojisym\triangle\mathrm{ABC}$ です。AB＝AD＋DB を先に求め、$\mathrm{AD}:\mathrm{AB}$ を相似比にします。
\end{minipage}\hfill
\begin{minipage}[t]{0.34\linewidth}\vspace{0pt}\centering
\begin{tikzpicture}[x=0.88mm, y=0.88mm]
  \draw[bkInk, line width=0.9pt] (14,27) -- (0,0) -- (36,0) -- cycle;
  \draw[bkOrange, line width=1.1pt] ($(14,27)!0.6667!(0,0)$) -- ($(14,27)!0.6667!(36,0)$);
  \node[font=\sffamily\footnotesize, anchor=south] at (14,27.6) {A};
  \node[font=\sffamily\footnotesize, anchor=north east] at (-0.4,-0.4) {B};
  \node[font=\sffamily\footnotesize, anchor=north west] at (36.4,-0.4) {C};
  \node[font=\sffamily\footnotesize, anchor=east] at ($(14,27)!0.6667!(0,0)+(-0.8,0)$) {D};
  \node[font=\sffamily\footnotesize, anchor=west] at ($(14,27)!0.6667!(36,0)+(0.8,0)$) {E};
\end{tikzpicture}
\end{minipage}\par
\shiki AB＝$4+2=6$。相似比は $\mathrm{AD}:\mathrm{AB}=4:6=2:3$\par
\hspace*{2.6\zw}$\mathrm{AE}:\mathrm{AC}=2:3$ より、$6:\mathrm{AC}=2:3$。$2\times\mathrm{AC}=18$、AC＝9\par
\hspace*{2.6\zw}EC＝AC－AE＝$9-6=3$\par
\hspace*{2.6\zw}確かめ：\mbox{AE : EC＝6 : 3＝2 : 1}、\mbox{AD : DB＝4 : 2＝2 : 1} で同じ\par
\hspace*{2.6\zw}$\mathrm{DE}=\mathrm{BC}\times\dfrac{2}{3}=9\times\dfrac{2}{3}=6$\par
\kotae EC＝3、DE＝6
\end{reidai}

\begin{jibun}
三角形ABCで、DE∥BC（Dは辺AB上、Eは辺AC上）です。AD＝3、DB＝6、AE＝2、BC＝12のとき、ECとDEの長さを求めましょう。
\bkjans{AB＝9、相似比 $3:9=1:3$。$2:\mathrm{AC}=1:3$ より AC＝6、EC＝$6-2=4$。$\mathrm{DE}=12\times\dfrac{1}{3}=4$。}
\end{jibun}
\bkivdpoint{平行線があったら、相似な2つの三角形を取り出して、ノートに別々にかき直すと、まちがえにくくなります。ABはAD＋DBです。}

\begin{reidai}[ピンホールカメラの像の大きさ]
\noindent\begin{minipage}[t]{0.56\linewidth}
小さな穴をあけた箱（ピンホールカメラ）で、穴から\qty{4}{m}はなれた所に立つ、身長\qty{1.6}{m}の人をうつします。穴から、像がうつる面までの深さは\qty{10}{cm}です。像の高さは何cmですか。\par
\kangaekata 光は穴を通ってまっすぐ進むので、X字形の2つの三角形ができます。相似比は、穴からの距離の比です。単位をcmにそろえます。
\end{minipage}\hfill
\begin{minipage}[t]{0.4\linewidth}\vspace{0pt}\centering
\begin{tikzpicture}[x=1mm, y=1mm]
  % 人（高さ 16）を左、穴 O を (24,8)、箱の奥を x=32
  \bkperson{(0,0)}{16}{bkBlue}
  \draw[bkOrange, line width=0.6pt] (0,16) -- (32,5.33);
  \draw[bkOrange, line width=0.6pt] (0,0) -- (32,10.67);
  \draw[bkInk!70, line width=1pt] (24,7.4) -- (24,-3) -- (32,-3) -- (32,19) -- (24,19) -- (24,8.6);
  \draw[bkBlue, line width=1.4pt] (31.4,10.67) -- (31.4,5.33);
  \node[font=\sffamily\footnotesize, anchor=north] at (12,-1.4) {\qty{4}{m}};
  \node[font=\sffamily\footnotesize, anchor=south] at (28,19.6) {\qty{10}{cm}};
\end{tikzpicture}
\end{minipage}\par
\shiki \qty{4}{m}＝\qty{400}{cm}、\qty{1.6}{m}＝\qty{160}{cm}\par
\hspace*{2.6\zw}相似比は $400:10=40:1$。像は人の $\dfrac{1}{40}$\par
\hspace*{2.6\zw}$160\times\dfrac{1}{40}=4$\par
\kotae \qty{4}{cm}（上下がさかさまの像）
\end{reidai}

\begin{jibun}
穴から\qty{10}{m}はなれた高さ\qty{5}{m}の木を、深さ\qty{8}{cm}のピンホールカメラでうつします。像の高さは何cmですか。
\bkjans{$\qty{1000}{cm}:\qty{8}{cm}=125:1$。$500\div125=\qty{4}{cm}$。}
\end{jibun}
\bkivdpoint{X字形は、上下がさかさまになった相似です。相似比は、交わる点（穴）からの距離の比です。}

\begin{reidai}[面積の比と体積の比]
（1）\nobreak 相似な2つの立体があり、相似比は $2:3$ です。小さいほうの表面積が\qty{36}{cm^2}、体積が\qty{24}{cm^3}のとき、大きいほうの表面積と体積を求めましょう。\par
（2）\nobreak 縮尺 $1:20$ の船の模型をぬるのに、ペンキが\qty{15}{mL}いりました。本物の船を同じ厚さでぬるには、ペンキは何Lいりますか。\par
\kangaekata 面積は $k^2$ 倍、体積は $k^3$ 倍です。ペンキの量は、ぬる面積に比例します。\par
\shiki （1）$k=\dfrac{3}{2}$。表面積は $36\times\left(\dfrac{3}{2}\right)^2=36\times\dfrac{9}{4}=81$\par
\hspace*{4.3\zw}体積は $24\times\left(\dfrac{3}{2}\right)^3=24\times\dfrac{27}{8}=81$\par
\hspace*{2.6\zw}（2）$k=20$。面積は $20^2=400$ 倍なので、$15\times400=\qty{6000}{mL}$\par
\hspace*{4.3\zw}$\qty{6000}{mL}\div\qty{1000}{mL/L}=\qty{6}{L}$\par
\kotae （1）\nobreak 表面積\qty{81}{cm^2}、体積\qty{81}{cm^3}\quad（2）\nobreak \qty{6}{L}
\end{reidai}

\begin{jibun}
（1）\nobreak 相似比 $3:4$ の2つの立体で、小さいほうの表面積が\qty{45}{cm^2}、体積が\qty{54}{cm^3}です。大きいほうの表面積と体積を求めましょう。（2）\nobreak 縮尺 $1:25$ の模型をぬるのに、ペンキが\qty{8}{mL}いりました。本物には何Lいりますか。
\bkjans{（1）\nobreak $45\times\dfrac{16}{9}=\qty{80}{cm^2}$、$54\times\dfrac{64}{27}=\qty{128}{cm^3}$。（2）\nobreak $8\times625=\qty{5000}{mL}$ で\qty{5}{L}。}
\end{jibun}
\bkivdpoint{長さが $k$ 倍なら、面積は $k^2$ 倍、体積は $k^3$ 倍です。まず $k$ を求めてから、2乗・3乗します。}

\begin{reidai}[すいの断面]
（1）\nobreak 高さ\qty{12}{cm}、底面の半径\qty{4}{cm}の円すいを、頂点から\qty{9}{cm}の所で底面に平行に切ります。断面の半径と面積を求めましょう。\par
（2）\nobreak 高さ\qty{10}{cm}、底面が1辺\qty{6}{cm}の正方形の四角すいを、頂点から\qty{5}{cm}の所で底面に平行に切ります。断面の1辺と面積を求めましょう。\par
\kangaekata 頂点から切りとった小さいすいは、もとのすいと相似で、相似比は「頂点からの長さ：高さ」です。\par
\shiki （1）相似比 $9:12=3:4$。半径は $4\times\dfrac{3}{4}=3$\par
\hspace*{4.3\zw}面積は $\pi\times3^2=9\pi\approx28.3$\par
\hspace*{2.6\zw}（2）相似比 $5:10=1:2$。1辺は $6\times\dfrac{1}{2}=3$、面積は $3\times3=9$\par
\hspace*{4.3\zw}確かめ：底面の面積36の $\left(\dfrac{1}{2}\right)^2=\dfrac{1}{4}$ 倍で9\par
\kotae （1）\nobreak 半径\qty{3}{cm}、面積 $9\pi$\,\unit{cm^2}（約\qty{28.3}{cm^2}）\quad（2）\nobreak 1辺\qty{3}{cm}、面積\qty{9}{cm^2}
\end{reidai}

\begin{jibun}
（1）\nobreak 高さ\qty{15}{cm}、底面の半径\qty{5}{cm}の円すいを、頂点から\qty{6}{cm}の所で底面に平行に切ります。断面の半径と面積を求めましょう。（2）\nobreak 高さ\qty{8}{cm}、底面が1辺\qty{4}{cm}の正方形の四角すいを、頂点から\qty{6}{cm}の所で底面に平行に切ります。断面の1辺と面積を求めましょう。
\bkjans{（1）\nobreak 相似比 $6:15=2:5$。半径 $5\times\dfrac{2}{5}=\qty{2}{cm}$、面積 $4\pi$\,\unit{cm^2}（約\qty{12.6}{cm^2}）。（2）\nobreak 相似比 $6:8=3:4$。1辺 $4\times\dfrac{3}{4}=\qty{3}{cm}$、面積\qty{9}{cm^2}。}
\end{jibun}
\bkivdpoint{頂点から切りとった小さいすいは、もとのすいと相似です。相似比は「頂点からの長さ：高さ」です。}

\begin{reidai}[直角三角形の辺の比を使う]
（1）\nobreak 長さ\qty{6}{m}のすべり台が、地面と30°の角をつくっています。すべり台のいちばん上の高さは何mですか。\par
（2）\nobreak 角が20°の直角三角形では、「たて÷横」が0.364です。角20°の坂を、横に\qty{50}{m}進むと、何m上がりますか。\par
\kangaekata （1）30°の直角三角形では、たての辺は斜辺の半分です。すべり台が斜辺です。（2）たて＝横×0.364 です。\par
\begin{bkcenter}
\begin{tikzpicture}[x=1mm, y=1mm]
  % （1）斜辺 6 m（30 mm）、30°。（2）横 50 m（40 mm）、20°
  \fill[bkBlue!12] (0,0) -- (25.98,0) -- (25.98,15) -- cycle;
  \draw[bkBlue, line width=1.4pt] (0,0) -- (25.98,15);
  \draw[bkInk, line width=0.6pt] (0,0) -- (25.98,0) -- (25.98,15);
  \bkivdright{(25.98,0)}{180}{90}
  \draw[bkInk, line width=0.5pt] (6,0) arc (0:30:6);
  \node[font=\sffamily\footnotesize] at (15:9) {30°};
  \node[font=\sffamily\footnotesize, anchor=south east] at (12.5,8) {\qty{6}{m}};
  \node[font=\sffamily\footnotesize, anchor=west] at (27,7.5) {？};
  \node[font=\sffamily\small\bfseries, anchor=east] at (-2,12) {（1）};
  \begin{scope}[shift={(56,0)}]
    \fill[bkGreen!15] (0,0) -- (40,0) -- (40,14.56) -- cycle;
    \draw[bkGreen!60!black, line width=1.4pt] (0,0) -- (40,14.56);
    \draw[bkInk, line width=0.6pt] (0,0) -- (40,0) -- (40,14.56);
    \bkivdright{(40,0)}{180}{90}
    \draw[bkInk, line width=0.5pt] (8,0) arc (0:20:8);
    \node[font=\sffamily\footnotesize] at (10:11) {20°};
    \node[font=\sffamily\footnotesize, anchor=north] at (20,-1.4) {横 \qty{50}{m}};
    \node[font=\sffamily\footnotesize, anchor=west] at (41,7.3) {？};
    \node[font=\sffamily\small\bfseries, anchor=east] at (-2,12) {（2）};
  \end{scope}
\end{tikzpicture}
\end{bkcenter}

\shiki （1）高さ＝$6\times\dfrac{1}{2}=3$\par
\hspace*{2.6\zw}（2）たて＝$50\times0.364=18.2$\par
\kotae （1）\nobreak \qty{3}{m}\quad（2）\nobreak 約\qty{18.2}{m}
\end{reidai}

\begin{jibun}
（1）\nobreak たこあげで、糸を\qty{8}{m}出したら、糸がまっすぐにはって、地面と30°の角をつくりました。たこは、手よりも何m高い所にありますか。（2）\nobreak 角20°の坂を、横に\qty{30}{m}進むと、約何m上がりますか（たて÷横＝0.364。小数第1位まで）。
\bkjans{（1）\nobreak $8\times\dfrac{1}{2}=4$ で、\qty{4}{m}高い。（2）\nobreak $30\times0.364=10.92$ で、約\qty{10.9}{m}。}
\end{jibun}
\bkivdpoint{30°の直角三角形は「正三角形の半分」と考えると、たての辺が斜辺の半分であることがすぐわかります。ほかの角では、比の値（0.364など）を使います。4-6では、この値を電卓で求めます。}

%% ---------------- まちがえやすいところ ----------------
\Needspace{14\baselineskip}
\bksection{tsumazuki}{まちがえやすいところをチェック}

\begin{tsumazuki}{対応する辺を、図の見た目で決めてしまう}
$\triangle\mathrm{ABC}\sojisym\triangle\mathrm{DEF}$ なのに、図が回っていたり、裏返っていたりすると、ちがう辺を組にしてしまいます。
\bkng{図で下にある辺どうし、BCとDEを組にする}
\bkok{記号の順番どおり、ABとDE、BCとEF、CAとFDを組にする}
相似の記号は、対応する頂点の順に書く約束です。迷ったら、等しい角どうしを先に見つけましょう。
\end{tsumazuki}

\begin{tsumazuki}{面積の比を、相似比と同じにする}
\bkng{相似比 $1:2$ なので、面積も2倍}
\bkok{面積は $2^2=4$ 倍。体積なら $2^3=8$ 倍}
たても横も2倍になるので、2を2回かけます。方眼の図を思い出そう。
\end{tsumazuki}

\begin{tsumazuki}{平行線と線分の比で、\mbox{DE : BC} を \mbox{AD : DB} にする}
\bkng{AD＝4、DB＝2 なので、\mbox{DE : BC＝4 : 2}}
\bkok{DE : BC＝\mbox{AD : AB＝4 : 6}}
相似な三角形は$\triangle\mathrm{ADE}$と$\triangle\mathrm{ABC}$です。BCに対応するのは、辺DBではなく、辺AB全体（AD＋DB）です。
\end{tsumazuki}

\begin{tsumazuki}{縮尺の計算で、単位を直しわすれる}
\bkng{$1:\num{25000}$ の地図で\qty{4}{cm}は、\num{100000}\,m}
\bkok{$4\times\num{25000}=\qty{100000}{cm}=\qty{1}{km}$}
地図の上の長さをcmではかったら、かけ算の答えもcmです。\num{100000}\,mは\qty{100}{km}で、ありえないほど長くなります。
\end{tsumazuki}

\begin{tsumazuki}{たてと横を別の倍率でのばしても、相似だと思う}
\bkng{写真を、横だけ2倍にのばしても、にた形}
\bkok{たても横も同じ倍率でのばしたときだけ、相似}
横だけのばすと、角の大きさが変わります。顔が横に太って見えるのは、そのためです。
\end{tsumazuki}

\begin{column}{怪獣は、立っていられる？}
身長\qty{1.7}{m}の人を、身長\qty{50}{m}まで、形を変えずに大きくしたとします。倍率は $k=50\div1.7\approx29.4$ です。体重は $k^3\approx\num{25400}$ 倍、足の骨の断面の面積は $k^2\approx865$ 倍。骨の断面1\,\unit{cm^2}あたりにかかる重さは、$k\approx29$ 倍になります。人の骨は、ふだんかかる重さの何倍かには耐えられますが、29倍ではおれてしまうでしょう。映画の巨大な怪獣が立っていられるのは、映画の中だけかもしれません。反対に、アリが体重の何倍もの物を運べるのも、小さい生き物ほど、体重にくらべて筋肉や骨の断面が大きいからです。
\end{column}

%% ---------------- 物理・工学 ----------------
\begin{butsuri}
\par\addvspace{3pt}\textbf{ピンホールカメラとレンズ}\par\noindent
\begin{minipage}[c]{0.62\linewidth}
ピンホールカメラでは、物の上のはしと下のはしから出た光が、小さな穴で交わり、X字形をつくります。2つの三角形は相似なので、像の大きさは「穴から像までの距離÷穴から物までの距離」倍になります。カメラのレンズでも、レンズの中心を通る光はまっすぐ進むので、同じX字形の相似で像の大きさが決まります。スマートフォンのカメラでは、レンズから像までが約\qty{4}{mm}なので、\qty{2}{m}先の身長\qty{1.6}{m}の人の像は、$1.6\times\dfrac{0.004}{2}=\qty{0.0032}{m}$、約\qty{3.2}{mm}です。
\end{minipage}\hfill
\begin{minipage}[c]{0.34\linewidth}\centering
\begin{tikzpicture}[x=0.78mm, y=0.78mm]
  % 凸レンズ：中心 (26,0)。物（矢印）は x=0 で高さ 12、像は x=40 で高さ −6.5
  \draw[bkGray, line width=0.5pt] (-2,0) -- (46,0);
  \draw[bkBlue!80!black, line width=0.7pt, fill=bkBlue!15] (26,0) ellipse (2.2 and 15);
  \draw[bk arrow, bkGreen!60!black, line width=1.4pt] (0,0) -- (0,12);
  \draw[bk arrow, bkGreen!60!black, line width=1.4pt] (40,0) -- (40,-6.46);
  \draw[bkOrange, line width=0.7pt] (0,12) -- (40,-6.46);
  \node[font=\sffamily\footnotesize, anchor=south] at (0,12.6) {物};
  \node[font=\sffamily\footnotesize, anchor=north] at (40,-7.2) {像};
  \node[font=\sffamily\footnotesize, anchor=south] at (26,15.6) {レンズ};
\end{tikzpicture}
\end{minipage}

\par\addvspace{4pt}\textbf{てこ}\par\noindent
支点から\qty{20}{cm}の所に荷物をのせ、支点から\qty{60}{cm}の所をおします。おす側を\qty{3}{cm}下げると、棒は支点のまわりに回り、X字形の相似な三角形ができます。相似比は $60:20=3:1$ なので、荷物は $3\div3=\qty{1}{cm}$ 上がります。動く長さが $\dfrac{1}{3}$ になるかわりに、荷物を持ち上げる力は、おす力の3倍になります（力と動く長さの関係は、物理シリーズで「仕事」を使ってくわしく調べます）。
\begin{bkcenter}
\begin{tikzpicture}[x=1mm, y=1mm]
  % 支点 (15,0)。左のうで 15 mm（20 cm）、右のうで 45 mm（60 cm）。棒の傾き −8/60
  \fill[bkGray!60] (15,0) -- (12,-6) -- (18,-6) -- cycle;
  \draw[bkGray, line width=0.4pt, dash pattern=on 1.5pt off 1.5pt] (-4,0) -- (64,0);
  \draw[bkStem, line width=2pt] (-4.5,{2+8*4.5/60}) -- (60,-6);
  % 荷物は板と同じ角度に傾け、下の辺を板の上の面にぴったりのせる（中心は支点から 15 mm＝20 cm）
  \begin{scope}[shift={(0,2)}, rotate={atan(-8/60)}]
    \filldraw[fill=bkOrange!70, draw=bkOrange!80!black, line width=0.4pt] (-3.5,0.35) rectangle (3.5,6.35);
  \end{scope}
  \node[font=\sffamily\footnotesize, anchor=east] at (-6,5.4) {荷物};
  \draw[bk arrow, bkBlue, line width=1.1pt] (60,5) -- (60,-4);
  \node[font=\sffamily\footnotesize, text=bkBlue!80!black, anchor=west] at (61,3) {おす};
  \node[font=\sffamily\footnotesize, anchor=south west] at (6,3.1) {\qty{20}{cm}};
  \node[font=\sffamily\footnotesize, anchor=south] at (37.5,0.4) {\qty{60}{cm}};
  \node[font=\sffamily\footnotesize, anchor=north] at (15,-6.8) {支点};
\end{tikzpicture}
\end{bkcenter}
\end{butsuri}

\begin{butsuri}
\par\addvspace{3pt}\textbf{縮尺模型と風洞}\par\noindent
\begin{minipage}[c]{0.6\linewidth}
新しい自動車や飛行機、高い建物を作るときは、縮尺 $1:5$ や $1:100$ などの模型を作り、トンネルのような装置（風洞）で風を当てて、空気の流れ方を調べます。縮尺 $1:5$ の模型では、風を受ける正面の面積は $\left(\dfrac{1}{5}\right)^2=\dfrac{1}{25}$ です。模型は、空気の流れ方まで本物とにるように、風の速さなどを調節して実験します。これも相似則の使い方です。
\end{minipage}\hfill
\begin{minipage}[c]{0.36\linewidth}\centering
\begin{tikzpicture}[x=1mm, y=1mm]
  \draw[bkInk!60, line width=0.8pt] (0,0) -- (46,0) (0,22) -- (46,22);
  \bkcar{(16,2)}{0.6}{bkRed}
  \foreach \yy in {6,11,16} {\draw[bkBlue, line width=0.7pt, -{Stealth[length=2.5pt]}] (1,\yy) -- (11,\yy);}
  \node[font=\sffamily\footnotesize, anchor=north] at (23,-1.4) {風洞の中の模型};
\end{tikzpicture}
\end{minipage}

\end{butsuri}

%% ---------------- 章末問題 ----------------
\bkendproblems
\bkproblemgroup{基本}{2}

\bkproblem 縮尺 $1:\num{50000}$ の地図があります。（1）\nobreak 地図の上で\qty{3}{cm}は、本物で何kmですか。（2）\nobreak 本物で\qty{4}{km}は、地図の上で何cmですか。（\bkreidairef{1}）
\bkpans{（1）\nobreak $3\times\num{50000}=\qty{150000}{cm}=\qty{1.5}{km}$。（2）\nobreak \qty{400000}{cm}÷\num{50000}＝\qty{8}{cm}。}

\bkproblem $\triangle\mathrm{ABC}\sojisym\triangle\mathrm{DEF}$ で、AB＝\qty{8}{cm}、BC＝\qty{10}{cm}、CA＝\qty{6}{cm}、DE＝\qty{12}{cm}です。相似比と、EF、FDを求めましょう。（\bkreidairef{2}）
\bkpans{相似比 $8:12=2:3$。$\mathrm{EF}=10\times\dfrac{3}{2}=\qty{15}{cm}$、$\mathrm{FD}=6\times\dfrac{3}{2}=\qty{9}{cm}$。}

\bkproblem 相似かどうかと、相似条件・相似比を答えましょう。（1）\nobreak 3辺が 2、3、4 と、5、7.5、10\quad（2）\nobreak 2つの角が35°、65°と、65°、80°（\bkreidairef{3}）
\bkpans{（1）\nobreak 相似。3辺の比、$2:5$。（2）\nobreak 3つめの角は80°と35°。どちらも35°、65°、80°なので相似（2つの角）。}

\bkproblem 長さ\qty{1.5}{m}の棒の影が\qty{2}{m}のとき、同じ時刻に影が\qty{12}{m}の電柱の高さを求めましょう。（\bkreidairef{4}）
\bkpans{高さ÷影＝$1.5\div2=0.75$。$12\times0.75=\qty{9}{m}$。}

\bkproblem 三角形ABCで、DE∥BC（D、Eは辺AB、AC上）、AD＝6、DB＝3、AE＝4、BC＝15のとき、ECとDEを求めましょう。（\bkreidairef{5}）
\bkpans{AB＝9、相似比 $6:9=2:3$。$4:\mathrm{AC}=2:3$ より AC＝6、EC＝2。$\mathrm{DE}=15\times\dfrac{2}{3}=10$。}

\bkproblem 相似比 $1:3$ の2つの立体で、小さいほうの表面積が\qty{20}{cm^2}、体積が\qty{8}{cm^3}です。大きいほうの表面積と体積を求めましょう。（\bkreidairef{7}）
\bkpans{表面積 $20\times9=\qty{180}{cm^2}$、体積 $8\times27=\qty{216}{cm^3}$。}

\bkproblem 高さ\qty{9}{cm}、底面の半径\qty{6}{cm}の円すいを、頂点から\qty{3}{cm}の所で底面に平行に切ります。断面の半径と面積を求めましょう。（\bkreidairef{8}）
\bkpans{相似比 $3:9=1:3$。半径 $6\times\dfrac{1}{3}=\qty{2}{cm}$、面積 $4\pi$\,\unit{cm^2}（約\qty{12.6}{cm^2}）。}

\bkproblemgroup{チャレンジ}{3}

\bkproblem 穴から\qty{20}{m}はなれた高さ\qty{12}{m}の建物を、深さ\qty{5}{cm}のピンホールカメラでうつします。像の高さは何cmですか。（\bkreidairef{6}）
\bkpans{$\qty{2000}{cm}:\qty{5}{cm}=400:1$。$1200\div400=\qty{3}{cm}$。}

\bkproblem 縮尺 $1:10$ の自動車の模型をぬるのに、ペンキが\qty{40}{mL}いりました。本物を同じ厚さでぬるには、何Lいりますか。（\bkreidairef{7}）
\bkpans{面積は $10^2=100$ 倍。$40\times100=\qty{4000}{mL}$ で\qty{4}{L}。}

\bkproblem （1）\nobreak 地面と30°、ななめの長さ\qty{12}{m}の駅のエスカレーターは、何m上がりますか。（2）\nobreak 角20°の坂（たて÷横＝0.364）で、横に\qty{80}{m}進むと約何m上がりますか（小数第1位まで）。（\bkreidairef{9}）
\bkpans{（1）\nobreak $12\times\dfrac{1}{2}=\qty{6}{m}$。（2）\nobreak $80\times0.364=29.12$ で、約\qty{29.1}{m}。}

\begin{furikaeri}
\begin{checklist}
\item 拡大図・縮図と縮尺を使って、本当の長さや図の上の長さを求められる。\bkfurref{bk:4-3:kakudai}
\item 三角形の相似条件を見つけ、相似比や平行線と線分の比から、辺の長さを求められる。\bkfurref{bk:4-3:jouken}
\item 面積の比は $k^2$、体積の比は $k^3$ を使い、すいの断面や直角三角形の辺の比を説明できる。\bkfurref{bk:4-3:menseki}
\end{checklist}
\end{furikaeri}

\bkfukidashi{delta}{形が同じなら、比も同じ。次は三平方の定理だよ。}
